\section{动力电池包装配分析}
\subsection{产品构成}
动力电池包并非单个电芯或简单模组，而是由电芯、模组、箱体、托盘、冷却板、绝缘件、母排、连接片、线束、采样单元、高压连接器以及 BMS 等部件共同构成的复杂机电系统。其产品特征决定了装配任务具有明显的跨学科耦合属性：结构装配需要保证尺寸链和刚度，电连接需要保证导通性能与接触电阻，热管理相关部件需要保证贴合与密封，控制部件安装则要求接口正确、线束走向规范，并满足后续测试需求 \cite{Kampker2021JRem,Harper2019Nature,Baazouzi2022ApEnergy}。

若从系统层级进一步拆解，动力电池包通常并不是“电芯直接进箱体”的简单堆叠，而是沿着“电芯--模组--电池包”的方式组织，并在包级同时集成冷却单元和电子控制单元。图~\ref{fig:product-architecture-pack}给出的模块化结构示意清楚展示了这一层级：模组内部既有电芯与接触系统，也包含局部从控 BMS；整包层则进一步叠加箱体、冷却和电子盒。这种层级关系意味着装配分析必须同时关注三个约束方向，即模组内部连接一致性、模组间集成关系以及整包级接口协调 \cite{Kampker2021JRem}。

\begin{figure}[H]
  \centering
  \includegraphics[width=0.74\textwidth]{fig13_product_architecture.png}
  \caption{动力电池包的模块化产品架构示意}
  \label{fig:product-architecture-pack}
\end{figure}

图~\ref{fig:product-pack-photo}给出了动力电池包在封装状态与开盖状态下的实物形态。封装状态对应运输、入线和整包装配后的外部观察对象，开盖状态则能够直接看到模组阵列、汇流排和内部布置关系。这种“外部构型 + 内部布局”的对照，对于理解后续模组入箱、线束接入、母排连接和终检位置选择具有直接帮助 \cite{Villagrossi2023JRem}。

\begin{figure}[H]
  \centering
  \includegraphics[width=0.64\textwidth]{fig6_battery_pack_photo.png}
  \caption{动力电池包封装状态与内部布局}
  \label{fig:product-pack-photo}
\end{figure}

除实物照片外，CAD 模型还适合呈现动力电池包在工艺规划阶段的关键作业与安全点位。图~\ref{fig:product-cad-model}展示了箱体模型上与拆装和检测密切相关的若干位置，如测温点、吊装点、排液点、排气口以及电量测量点。虽然该图直接服务于拆解场景，但对装配分析同样具有参考价值，因为吊装路径、检测位置和安全接口布置本身就是装配规划、工位设计与后续维护性分析的重要约束 \cite{Villagrossi2023JRem}。这也说明，动力电池包装配研究并不只是面向实物执行，还需要在数字模型层面对结构关系、风险位置和作业边界进行预先理解。

\begin{figure}[H]
  \centering
  \includegraphics[width=0.76\textwidth]{fig_cad_model_fiat_500e_lr.png}
  \caption{动力电池包 CAD 模型中的关键作业与安全点位}
  \label{fig:product-cad-model}
\end{figure}

\subsection{装配任务与场景特征}
从任务分解角度看，动力电池包的制造可分为模组级任务和整包级任务两类。模组级任务主要包括电芯编组、极耳或连接片连接、模组框架固定、绝缘隔离与局部检测等；整包级任务则进一步延伸到箱体准备、模组入箱、冷却系统集成、母排连接、焊接、线束与 BMS 安装、密封、功能测试与终检。与一般装配过程不同，动力电池包在多个工位均伴随检测动作，例如入箱前的尺寸与外观检查、焊接后的电阻与外观检查、封装前的位置与间隙检查，以及终装后的绝缘、密封和功能测试。因此，装配与检测并非先后独立关系，而是天然交织的复合流程 \cite{Zhang2022JMSY,Wegener2022Energies,Mayoral2023IJPR}。

将动力电池包作为本文对象，还有一个重要原因，即其既具有足够的复杂性，也具有较强的典型性。其复杂性体现在同时涉及机械装配、电气连接、热管理集成和安全检测；其典型性则体现在代表了当前新能源装备制造中“高安全、高一致性、高柔性”的共性要求。相较于仅讨论单个零件或单道工序，以动力电池包为对象更能体现智能装配与检测的系统性，也更便于将机器人、机器视觉、人机协作、知识图谱和大模型智能体等技术置于同一条制造主线下展开讨论。

动力电池包的装配任务还具有显著的安全属性。一方面，装配过程中的焊接质量、绝缘状态、紧固力矩和密封质量直接影响产品服役可靠性；另一方面，工位上的误操作还可能引发局部短路、部件变形或返修成本上升。因此，动力电池包装配不宜被理解为单纯的“搬运和连接”，而应被视为贯穿结构、电气、热管理和质量保证的系统工程 \cite{Mayoral2023IJPR,Li2020TIE}。
